% GENERATED FILE. Edit canonical lessons, then run npm run generate:books.
% canonical-source-hash: fnv1a64:292cd03b867a8f3d
% canonical-lessons: MR-C408-navadi, MR-C408-valhe, MR-C408-valhavne, MR-C408-helkavne, MR-C408-nangar

\chapter{Boats and Oars}
\label{ch:mr-boats-and-oars}
\hlchaptermodality{408}

\begin{chapteropening}
\textbf{By the end of this chapter:} \emph{I can talk about a boatman, oars, rowing and a boat rocking at anchor.}

Complete the last lesson of chapter 408: I can talk about a boatman, oars, rowing and a boat rocking at anchor.
\end{chapteropening}

\section[\texorpdfstring{\mr{नावाडी} (nāvāḍī)}{नावाडी (nāvāḍī)}]{\mr{नावाडी} (nāvāḍī) — a boatman}

\label{lesson:MR-C408-navadi}



\begin{quote}
Before the new one: say the Marathi for a signpost, then the Marathi for a long trudge.
\end{quote}

\subsection*{You'll want to know: \mr{नावाडी}}
\textbf{\mr{नावाडी}} — \emph{nāvāḍī} — \textquotedblleft{}a boatman\textquotedblright{}.

\begin{quote}
\mr{नावाड्याने} \mr{आम्हाला} \mr{नदीपलीकडे} \mr{नेलं}.
\end{quote}

\begin{quote}
\emph{nāvāḍyāne āmhālā nadīpalīkaḍe nelaṁ.}
\end{quote}

\begin{quote}
The boatman took us across the river.
\end{quote}

Masculine.

Built on \mr{नाव}, a boat; before an ending it becomes \mr{नावाड्या}-.

\begin{cousinweb}
\textbf{\mr{नावाडी}} (\emph{nāvāḍī}) is inherited: CDIAL lists Marathi \emph{nāvāḍī}, \textquotedblleft{}boatman\textquotedblright{}, under Sanskrit \emph{nāvā}, \textquotedblleft{}boat\textquotedblright{} (CDIAL 7081).
\end{cousinweb}

\begin{grammarlens}[title={What you've built}]
One more word for getting around.
\end{grammarlens}

\subsection*{Your turn}
\begin{itemize}
\raggedright
  \item \emph{Say it:} \emph{nāvāḍī}
  \item \emph{Say it:} \emph{nāvāḍī}, once more
  \item \emph{Say it:} \emph{nāvāḍyāne āmhālā nadīpalīkaḍe nelaṁ.} — The boatman took us across the river.
  \item \emph{Recall:} say the Marathi for a signpost, then the Marathi for a long trudge, then say \emph{nāvāḍī} again
\end{itemize}

\begin{tcolorbox}[breakable,colback=teal!4,colframe=teal!35!black,title={Before you move on}]
What does \mr{नावाडी} mean? (\textquotedblleft{}A boatman\textquotedblright{}.) Say it once more.
\end{tcolorbox}
\section[\texorpdfstring{\mr{वल्हे} (valhe)}{वल्हे (valhe)}]{\mr{वल्हे} (valhe) — an oar}

\label{lesson:MR-C408-valhe}



\begin{quote}
Before the new one: say the Marathi for a long trudge, then the Marathi for a boatman.
\end{quote}

\subsection*{You'll want to know: \mr{वल्हे}}
\textbf{\mr{वल्हे}} — \emph{valhe} — \textquotedblleft{}an oar\textquotedblright{}.

\begin{quote}
\mr{होडीत} \mr{दोन} \mr{वल्ही} \mr{होती}.
\end{quote}

\begin{quote}
\emph{hoḍīt don valhī hotī.}
\end{quote}

\begin{quote}
There were two oars in the boat.
\end{quote}

Neuter. Plural: \mr{वल्ही}, \emph{valhī}.

A neuter noun: one \mr{वल्हे}, two \mr{वल्ही}.

\begin{cousinweb}
\textbf{\mr{वल्हे}} (\emph{valhe}) is a Marathi word for an oar; its origin is not settled.
\end{cousinweb}

\begin{grammarlens}[title={What you've built}]
One more word for getting around.
\end{grammarlens}

\subsection*{Your turn}
\begin{itemize}
\raggedright
  \item \emph{Say it:} \emph{valhe}
  \item \emph{Say it:} \emph{valhe}, once more
  \item \emph{Say it:} \emph{hoḍīt don valhī hotī.} — There were two oars in the boat.
  \item \emph{Recall:} say the Marathi for a long trudge, then the Marathi for a boatman, then say \emph{valhe} again
\end{itemize}

\begin{tcolorbox}[breakable,colback=teal!4,colframe=teal!35!black,title={Before you move on}]
What does \mr{वल्हे} mean? (\textquotedblleft{}An oar\textquotedblright{}.) Say it once more.
\end{tcolorbox}
\section[\texorpdfstring{\mr{वल्हवणे} (valhavṇe)}{वल्हवणे (valhavṇe)}]{\mr{वल्हवणे} (valhavṇe) — to row}

\label{lesson:MR-C408-valhavne}



\begin{quote}
Before the new one: say the Marathi for a boatman, then the Marathi for an oar.
\end{quote}

\subsection*{You'll want to know: \mr{वल्हवणे}}
\textbf{\mr{वल्हवणे}} — \emph{valhavṇe} — \textquotedblleft{}to row\textquotedblright{}.

\begin{quote}
\mr{नावाडी} \mr{हळूहळू} \mr{होडी} \mr{वल्हवत} \mr{होता}.
\end{quote}

\begin{quote}
\emph{nāvāḍī haḷūhaḷū hoḍī valhavat hotā.}
\end{quote}

\begin{quote}
The boatman was slowly rowing the boat.
\end{quote}

Past, with -\mr{ने} (he …): \mr{त्याने} \mr{वल्हवलं}, \emph{tyāne valhavlaṁ}.

Built on \mr{वल्हे}, an oar: \mr{होडी} \mr{वल्हवणे}, to row a boat.

\begin{cousinweb}
\textbf{\mr{वल्हवणे}} (\emph{valhavṇe}) is a Marathi verb built on \textbf{\mr{वल्हे}}, \textquotedblleft{}an oar\textquotedblright{}, with the ending -\mr{वणे}; the origin of the noun is not settled.
\end{cousinweb}

\begin{grammarlens}[title={What you've built}]
One more word for getting around.
\end{grammarlens}

\subsection*{Your turn}
\begin{itemize}
\raggedright
  \item \emph{Say it:} \emph{valhavṇe}
  \item \emph{Say it:} \emph{valhavṇe}, once more
  \item \emph{Say it:} \emph{nāvāḍī haḷūhaḷū hoḍī valhavat hotā.} — The boatman was slowly rowing the boat.
  \item \emph{Recall:} say the Marathi for a boatman, then the Marathi for an oar, then say \emph{valhavṇe} again
\end{itemize}

\begin{tcolorbox}[breakable,colback=teal!4,colframe=teal!35!black,title={Before you move on}]
What does \mr{वल्हवणे} mean? (\textquotedblleft{}To row\textquotedblright{}.) Say it once more.
\end{tcolorbox}
\section[\texorpdfstring{\mr{हेलकावणे} (helkāvṇe)}{हेलकावणे (helkāvṇe)}]{\mr{हेलकावणे} (helkāvṇe) — to toss about; to rock}

\label{lesson:MR-C408-helkavne}



\begin{quote}
Before the new one: say the Marathi for an oar, then the Marathi for to row.
\end{quote}

\subsection*{You'll want to know: \mr{हेलकावणे}}
\textbf{\mr{हेलकावणे}} — \emph{helkāvṇe} — \textquotedblleft{}to toss about; to rock\textquotedblright{}.

\begin{quote}
\mr{वादळात} \mr{होडी} \mr{हेलकावत} \mr{होती}.
\end{quote}

\begin{quote}
\emph{vādaḷāt hoḍī helkāvat hotī.}
\end{quote}

\begin{quote}
The boat was tossing about in the storm.
\end{quote}

Past (he …): \mr{तो} \mr{हेलकावला}, \emph{to helkāvlā}.

The noun \mr{हेलकावा} is one toss: \mr{हेलकावे} \mr{खाणे}, to be tossed about.

\begin{cousinweb}
\textbf{\mr{हेलकावणे}} (\emph{helkāvṇe}) is a Marathi verb for a boat or a load swinging from side to side; its origin is not settled.
\end{cousinweb}

\begin{grammarlens}[title={What you've built}]
One more word for getting around.
\end{grammarlens}

\subsection*{Your turn}
\begin{itemize}
\raggedright
  \item \emph{Say it:} \emph{helkāvṇe}
  \item \emph{Say it:} \emph{helkāvṇe}, once more
  \item \emph{Say it:} \emph{vādaḷāt hoḍī helkāvat hotī.} — The boat was tossing about in the storm.
  \item \emph{Recall:} say the Marathi for an oar, then the Marathi for to row, then say \emph{helkāvṇe} again
\end{itemize}

\begin{tcolorbox}[breakable,colback=teal!4,colframe=teal!35!black,title={Before you move on}]
What does \mr{हेलकावणे} mean? (\textquotedblleft{}To toss about; to rock\textquotedblright{}.) Say it once more.
\end{tcolorbox}
\section[\texorpdfstring{\mr{नांगर} (nāṅgar)}{नांगर (nāṅgar)}]{\mr{नांगर} (nāṅgar) — an anchor; a plough}

\label{lesson:MR-C408-nangar}



\begin{quote}
Before the new one: say the Marathi for to row, then the Marathi for to toss about.
\end{quote}

\subsection*{You'll want to know: \mr{नांगर}}
\textbf{\mr{नांगर}} — \emph{nāṅgar} — \textquotedblleft{}an anchor; a plough\textquotedblright{}.

\begin{quote}
\mr{जहाजाने} \mr{बंदरात} \mr{नांगर} \mr{टाकला}.
\end{quote}

\begin{quote}
\emph{jahājāne bandarāt nāṅgar ṭāklā.}
\end{quote}

\begin{quote}
The ship dropped anchor in the harbour.
\end{quote}

Masculine.

\mr{नांगर} \mr{टाकणे}, to drop anchor; on a farm the same word is the plough.

\begin{cousinweb}
\textbf{\mr{नांगर}} (\emph{nāṅgar}) is inherited in the sense \textquotedblleft{}plough\textquotedblright{}: CDIAL lists Marathi \emph{nā̃gar} under Sanskrit \emph{lāṅgala}, \textquotedblleft{}plough\textquotedblright{} (CDIAL 11006); how the word also came to mean a ship's anchor is not settled.
\end{cousinweb}

\begin{grammarlens}[title={What you've built}]
One more word for getting around.
\end{grammarlens}

\subsection*{Your turn}
\begin{itemize}
\raggedright
  \item \emph{Say it:} \emph{nāṅgar}
  \item \emph{Say it:} \emph{nāṅgar}, once more
  \item \emph{Say it:} \emph{jahājāne bandarāt nāṅgar ṭāklā.} — The ship dropped anchor in the harbour.
  \item \emph{Recall:} say the Marathi for to row, then the Marathi for to toss about, then say \emph{nāṅgar} again
\end{itemize}

\begin{tcolorbox}[breakable,colback=teal!4,colframe=teal!35!black,title={Before you move on}]
What does \mr{नांगर} mean? (\textquotedblleft{}An anchor; a plough\textquotedblright{}.) Say it once more.
\end{tcolorbox}
